\subsection{环上的矩阵}

数学中最重要的矩阵是环 A 上的矩阵。对应于指标集 I, K 的 A 上矩阵的集合 \(A^{I \times K}\) 于是典范地具有一个 (A, A)-双模结构（§1，第 14 号）。

对每个有序对 \((i,k)\in\mathbf{I}\times\mathbf{K}\) ，设 \(E_{ik}\) 是矩阵 \((a_{jl})\) ，它满足 \(a_{ik}=1\) 且 \(a_{jl}=0\) （对 \((j,l)\neq(i,k)\) ）；诸 \(E_{ik}\) 称为矩阵集合 \(\mathbf{A}^{\mathbf{I}\times\mathbf{K}}\) 中的矩阵单位；若 I 与 K 有限，它们构成该集合关于其左或右 A-模结构的典范基（§1，第 11 号）。显然

\[
{ } ^ { t } E _ { i k } = E _ { k i } .
\]

除非另有说明，A 上两个矩阵（假定有定义）的乘积 \(M'M''\) 将总是理解为相对于 A 中乘法 \((x,y)\mapsto xy\) 的乘积（或者说，“在 A 中计算”）。于是我们有（第 2 号）结合律与分配律公式

\[
(X Y) Z = X (Y Z)\tag{7}
\]

\[
\left\{ \begin{array}{l} X (Y + Z) = X Y + X Z \\ (X + Y) Z = X Z + Y Z \end{array} \right.\tag{8}
\]

对于 A 上的三个矩阵 X, Y, Z，只要这些公式中出现的和与积有定义。

特别是，如果 \(E_{ik}\) （相应地 \(E_{kl}'\) , \(E_{il}''\) ）分别是 \(\mathbf{A}^{\mathbf{I} \times \mathbf{K}}\) （相应地 \(\mathbf{A}^{\mathbf{K} \times \mathbf{L}}\) , \(\mathbf{A}^{\mathbf{I} \times \mathbf{L}}\) ）中的矩阵单位，且 \(\mathbf{I} = [1, p]\) , \(\mathbf{K} = [1, q]\) , \(\mathbf{L} = [1, r]\) ，我们得到公式

\[
\left\{ \begin{array}{l l} E _ {i k} E _ {j l} ^ {\prime} = 0 & \text { if } k \neq j \\ E _ {i k} E _ {k l} ^ {\prime} = E _ {i l} ^ {\prime \prime}. \end{array} \right.\tag{9}
\]

设 \(\mathbf{A}^0\) 是 \(\mathbf{A}\) 的反环，并用 \(a * b (= ba)\) 表示 \(a\) 与 \(b\) 在 \(\mathbf{A}^0\) 中的乘积；那么，对两个矩阵 \(X, Y\) （ \(\mathbf{A}\) 上的），其乘积有定义，

\[
{ } ^ { t } ( X Y ) = { } ^ { t } Y * { } ^ { t } X\tag{10}
\]

其中右边的 \(^{t}Y\) 与 \(^{t}X\) 被看作元素属于 \(A^{0}\) 的矩阵；当 A 交换时，则

\[
{ } ^ { t } ( X Y ) = { } ^ { t } Y . { } ^ { t } X\tag{11}
\]

命题 1. 设 A、B 是两个环，且 \(M = (m_{ik})_{(i,k) \in I \times K}\) 和

\[
M ^ {\prime} = \left(m _ {i k} ^ {\prime}\right) _ {(i, k) \in \mathbf {I} \times \mathbf {K}}
\]

两个具有有限指标集的矩阵，取值于一个 (A, B)-双模 G 上。假设对每个单行、元素属于 A 的矩阵单位 \(L = (a_{i})_{i \in I}\) 以及每个单列、元素属于 B 的矩阵单位 \(C = (b_{k})_{k \in K}\) ，都有 L.M.C = L.M'.C（乘积经由 (A, B)-模 G 的外运算计算）；则 \(M = M'\) 。

若把 \(L\) 取为矩阵单位 \((a_s)\) ，其中 \(a_i = 1\) , \(a_s = 0\) （对 \(s \neq i\) ），并把 \(C\) 取为矩阵单位 \((b_t)\) ，其中 \(b_k = 1\) , \(b_t = 0\) （对 \(t \neq k\) ），则乘积 \(L.M.C\) 与 \(L.M'.C\) 都是单元素矩阵，其元素分别等于 \(m_{ik}\) 与 \(m'_{ik}\) 。

设 A、B 是两个环， \(\sigma: A \to B\) 是一个同态。

对每个矩阵 \(M = (m_{\iota\kappa})\) （ \(\mathbf{A}\) 上的），我们将用 \(\sigma(M)\) 表示矩阵 \((\sigma(m_{\iota\kappa}))\) （ \(\mathbf{B}\) 上的）；显然 \(\sigma(aM) = \sigma(a)\sigma(M)\) , \(\sigma(Ma) = \sigma(M)\sigma(a)\) （对 \(a \in \mathbf{A}\) ），同时 \(\sigma(^t M) = ^t(\sigma(M))\) ，并且

\[
\left\{ \begin{array}{c} \sigma (M + M ^ {\prime}) = \sigma (M) + \sigma (M ^ {\prime}) \\ \sigma (M M ^ {\prime}) = \sigma (M) \sigma (M ^ {\prime}) \end{array} \right.\tag{12}
\]

当所考虑的运算有定义时，(12) 左边与右边的乘积分别在 A 与 B 中计算。当 \(\sigma\) 记作 \(x \mapsto x^{\sigma}\) 时，我们用 \(M^{\sigma}\) 代替 \(\sigma(M)\) 。

特别地，考虑 A 的一个反自同态 \(\sigma\) ，即 A 到反环 \(\mathbf{A}^0\) 的一个同态，或 A 到自身的一个映射，使得

\[
\sigma (a + a ^ {\prime}) = \sigma (a) + \sigma (a ^ {\prime}), \quad \sigma (a a ^ {\prime}) = \sigma (a ^ {\prime}) \sigma (a)
\]

对所有 \(a, a'\) （在 \(\mathbf{A}\) 中）；那么，对两个矩阵 \(M, M'\) （ \(\mathbf{A}\) 上的），其乘积 \(MM'\) 有定义时，

\[
\sigma (M M ^ {\prime}) = ^ {t} (\sigma (^ {t} M ^ {\prime}). \sigma (^ {t} M))\tag{13}
\]

其中两边的乘积均在A中计算；这由(10)和(12)直接得出。

